\chapter{Estado del arte}
\label{cap:estado-arte}

\section{Criterios de revisión}
La revisión parte de las referencias usadas durante el desarrollo y las organiza en tres grupos: redes inalámbricas de sensores, sistemas de adquisición y antecedentes de monitoreo de condición. Se consultaron publicaciones originales y documentación de fabricantes para separar los resultados experimentales de las especificaciones. Es una revisión focalizada en las decisiones del prototipo, no una búsqueda sistemática de todas las soluciones existentes.

Los sistemas se comparan por la variable medida, el medio de comunicación, dónde se almacenan los datos, qué archivos de diseño se publican y qué validación reportan. Apertura y operación local son propiedades distintas: un sistema puede usar un protocolo abierto con software comercial, o guardar los datos localmente con radios propietarios.

\section{Redes de sensores y comunicaciones de baja potencia}
Las revisiones de \citet{akyildiz2002} y \citet{yick2008} muestran que una red de sensores no se diseña al margen de sus restricciones de energía, procesamiento y despliegue, y que la adquisición y el transporte de la información se tratan juntos. Esa es su pertinencia para SmartSense.

\citet{centenaro2016} presentan las comunicaciones de baja tasa y largo alcance para Internet de las cosas; \citet{augustin2016} estudian el funcionamiento y el desempeño de LoRa y LoRaWAN con pruebas y simulaciones, y \citet{bor2016} examinan sus límites de escalabilidad. Estos trabajos respaldan la elección de la radio para mensajes breves y periódicos y muestran que la carga del canal, la configuración y el escenario forman parte de la evaluación. Sus cifras no se trasladan al enlace de SmartSense, que usa su propio formato de aplicación y pocos nodos.

\section{Sistemas de adquisición relacionados}
\citet{polonelli2019} desarrollan una red LoRaWAN alimentada por batería para observar temperatura y humedad en un centro de datos, con seis meses de operación y más de seis millones de puntos de datos; el gateway puede integrar el servidor y operar localmente. Su aporte es la evaluación prolongada de una red térmica, en un escenario ambiental distinto del contacto con la carcasa de una máquina.

HIGROTERM, de \citet{higroterm2021}, es un registrador abierto de laboratorio con sensores DHT22 cableados, pantalla y tarjeta SD, que publica componentes, código y documentación de construcción: un antecedente directo de adquisición local reproducible, sin red inalámbrica ni sondas RTD sobre equipos industriales. \citet{mois2017} analizan sensores inalámbricos para monitoreo ambiental, también con una variable y un escenario distintos de la temperatura superficial.

\citet{jakobsen2024} presenta un sensor abierto de monitoreo de condición con acelerómetros, micrófono y sensor de temperatura, que combina BLE para los datos y LoRaWAN para las características procesadas, evaluado con excitación mecánica y un banco con un rodamiento defectuoso. Es el antecedente más cercano en maquinaria; SmartSense se concentra, en cambio, en el historial térmico de una PT100 externa y en una interfaz local de consulta.

\section{Comparación normalizada}
La Tabla~\ref{tab:estado-arte-comparacion} resume solo las características respaldadas por las fuentes citadas. ``No especificado'' significa que la documentación examinada no lo dice, no que el sistema carezca de la propiedad, y la columna de limitación se refiere al problema de este trabajo, no al valor del antecedente.

\begingroup
\small
\setstretch{1.05}
\setlength{\tabcolsep}{3pt}
\begin{longtable}{@{}>{\raggedright\arraybackslash}p{2.1cm}>{\raggedright\arraybackslash}p{2cm}>{\raggedright\arraybackslash}p{2.1cm}>{\raggedright\arraybackslash}p{2.4cm}>{\raggedright\arraybackslash}p{1.7cm}>{\raggedright\arraybackslash}p{2.2cm}>{\raggedright\arraybackslash}p{2.25cm}@{}}
\caption{Comparación de sistemas de adquisición y monitoreo relacionados.}\label{tab:estado-arte-comparacion}\\
\toprule
Sistema & Sensor / variable & Comunicación & Almacenamiento local / nube & Apertura & Validación & Limitación \\
\midrule
\endfirsthead
\multicolumn{7}{l}{Tabla~\thetable{} (continuación)}\\
\toprule
Sistema & Sensor / variable & Comunicación & Almacenamiento local / nube & Apertura & Validación & Limitación \\
\midrule
\endhead
\bottomrule
\endfoot
Polonelli et al. (2019) & Temperatura y humedad ambiental & LoRaWAN & Gateway con opción de servidor local & Diseño descrito; licencia integral no especificada & Seis meses en centro de datos & No evalúa contacto térmico en motores. \\
\addlinespace
HIGROTERM (2021) & DHT22: temperatura y humedad & Sensores cableados & SD local & Código y diseño publicados & Ensayos de laboratorio & No incorpora enlace inalámbrico RTD. \\
\addlinespace
Jakobsen (2024) & Vibración, sonido y temperatura & BLE y LoRaWAN & Datos hacia receptor; historial local/nube no especificado & Archivos de diseño publicados & Excitación mecánica y rodamiento defectuoso & Su evaluación se centra en señales mecánicas. \\
\addlinespace
WISE-2410 & Acelerómetro triaxial y temperatura & LoRaWAN & Servidor de aplicación según infraestructura & Producto comercial; protocolo estándar & Especificación del fabricante & Historial y acceso dependen de la integración. \\
\addlinespace
SmartSense & PT100 externa: temperatura superficial & LoRa punto a punto & microSD / LittleFS; dashboard local & Diseño, firmware y software del proyecto & Piloto: \NumeroNodos{} nodos, \DuracionComunDias{} días & Concordancia de campo; energía y PDR estimados. \\
\end{longtable}
\endgroup

\section{Comparación con WISE-2410}
El WISE-2410 es la referencia comercial pertinente por su orientación al monitoreo de maquinaria: su ficha \citep{wise2410} especifica vibración triaxial, cálculo de indicadores, sensor de temperatura y carcasa IP66. La Tabla~\ref{tab:wise-smartsense} contrasta esas prestaciones con el alcance documentado del prototipo, característica por característica, e indica en cada fila quién queda mejor. Es una comparación de arquitectura y funciones; no hubo un ensayo controlado de los dos sistemas.

\begin{table}[htbp]
\centering
\caption{Comparación técnica de SmartSense y WISE-2410.}
\label{tab:wise-smartsense}
\small
\begin{tabularx}{\textwidth}{@{}>{\raggedright\arraybackslash}p{2.6cm}XX>{\raggedright\arraybackslash}p{2.4cm}@{}}
\toprule
Criterio & SmartSense & WISE-2410 & Ventaja \\
\midrule
Variable y montaje & PT100 externa de tres hilos en superficie. & Temperatura integrada y vibración triaxial. & SmartSense en contacto; WISE en vibración \\
Comunicación & LoRa punto a punto; gateway propio. & LoRaWAN; gateway y servidor de aplicación. & Distinta; sin ventaja \\
Historial y consulta & Archivos locales y dashboard modificable. & Funciones dependientes de la plataforma integrada. & SmartSense \\
Apertura & Fuentes, conexiones y diseño disponibles en el proyecto. & Producto comercial con interfaces documentadas. & SmartSense \\
Protección física & Prototipo; sin clasificación IP demostrada. & Carcasa IP66 especificada. & WISE \\
Energía & Autonomía estimada desde \DuracionDescarga{} de descarga parcial; periodo nominal \PeriodoNominal{}. & Hasta dos años especificados con actualización una vez por hora. & WISE por especificación; SmartSense pendiente \\
Alcance experimental & Piloto térmico y pruebas complementarias. & Prestaciones de ficha; comparación de planta limitada a la configuración disponible. & No comparable \\
\bottomrule
\end{tabularx}
\end{table}

La autonomía que declara el fabricante usa un intervalo de actualización distinto del ensayado en SmartSense; compararlas exigiría igualar carga de trabajo, condiciones ambientales y criterio de fin de descarga. Tener un dashboard propio facilita adaptar la consulta, aunque las plataformas que integran el WISE-2410 también pueden ofrecer históricos locales.

La BOM (\texttt{hardware/BOM.csv}), revisada el 15 de septiembre de 2026, suma 627\,000~COP para tres nodos, gateway y cableado compartido: 149\,000~COP por nodo en piezas específicas, 130\,000~COP de gateway y 50\,000~COP compartidos. Equivale aproximadamente a 157~USD en total y 37~USD por nodo con el supuesto de conversión del proyecto de 4000~COP/USD. Los precios son los de compra del proyecto, sin fecha ni tasa de cambio verificada, y no incluyen mano de obra, envío ni certificación; como no hay una cotización fechada del WISE-2410, no se afirma una ventaja económica cuantificada. Los nodos usan una celda LiPo de 1000~mAh, la misma capacidad del ensayo energético.

\section{Síntesis y contribución del trabajo}
Los antecedentes muestran soluciones abiertas y comerciales, con almacenamiento local o con servidor. La oportunidad que aborda este trabajo es integrar una sonda RTD externa, radio punto a punto, alimentación por batería e histórico local en un conjunto reproducible y ajustado al piloto. Ningún componente es nuevo por separado; el aporte está en la integración documentada y evaluada.

La literatura de mantenimiento va más allá de adquirir datos: procesamiento, diagnóstico y pronóstico \citep{jardine2006,cinar2020}. SmartSense se ubica en la adquisición y presentación de la información térmica, y así debe juzgarse: por la cadena funcional, la utilidad del historial y los límites declarados de su evidencia.
